% !TEX program = xelatex
% AI-effects 프로젝트 --- AEI V6(2026년 4--5월) Claude.ai 대화 수와 1P API 호출 수의
% 절대 규모를 외부자료(Similarweb 트래픽, Anthropic 공식 run-rate 매출)로 추정하는
% 방법을 정리한 작업 노트(report/). 파라미터 원자료: data/raw/macro/scaling/
% 컴파일: xelatex -> biber -> xelatex -> xelatex (kotex, biblatex/biber 필요)
\documentclass[11pt]{article}

\usepackage{fontspec}
\usepackage{kotex}
\usepackage[a4paper,margin=2.2cm]{geometry}
\usepackage{longtable}
\usepackage{booktabs}
\usepackage{array}
\usepackage{amsmath}
\usepackage{amssymb}
\usepackage{xcolor}
\usepackage{hyperref}
\usepackage{enumitem}
\usepackage{titlesec}
\usepackage[backend=biber,style=authoryear,maxcitenames=2,uniquelist=false,sorting=nyt]{biblatex}
\addbibresource{../../reference/references.bib}

\hypersetup{colorlinks=true,linkcolor=blue!50!black,citecolor=blue!50!black,urlcolor=blue!50!black}
\newcolumntype{L}[1]{>{\raggedright\arraybackslash}p{#1}}
\setlength{\LTleft}{0pt}
\setlength{\LTright}{0pt}
\renewcommand{\arraystretch}{1.15}
\setlength{\tabcolsep}{4pt}

\title{AEI 사용량의 절대 규모를\\[1mm]
어떻게 추정할 것인가\\[2mm]
\large --- V6(2026년 4--5월) Claude.ai 대화 수와 1P API 호출 수의 외부자료 기반 추정 방법 ---}
\author{AI-effects 프로젝트}
\date{\today}

\begin{document}
\maketitle

\begin{abstract}
\noindent AEI V6(\texttt{release\_2026\_06\_26})는 플랫폼 안의 점유율(\texttt{pct})과 과업별 절감시간만 공개하고, 대화$\cdot$호출 \textbf{건수}는 공개하지 않는다. 그래서 노동비용등가(LCE)처럼 수준(level)이 필요한 지표를 계산하려면 사용량의 절대 규모를 외부 자료로 따로 추정해야 한다. 본 노트는 \textcite{fan-2026-aggregate-gains-from-ai}가 쓴 규모 가정을 원문에서 다시 정리하고, 이를 V6 기간(2026년 4--5월)으로 옮길 수 있는 외부 자료와 추정식을 플랫폼별로 제시한다. 결론은 다음과 같다. (1) \textbf{Claude.ai}는 Fan and Nguyen의 방식(2026년 2월 주 2억 건을 Similarweb 월 방문 수 비율로 늘리기)을 그대로 이으면 2026년 4월 주 약 5.7억 건, 5월 주 약 6.6억 건이 된다. 다만 앱 사용의 급증, Cowork 포함, 시장점유율 변화 때문에 이 값은 기준값일 뿐이며 민감도 분석이 반드시 필요하다. (2) \textbf{1P API}는 Anthropic이 공식 발표한 run-rate 매출에서 호출 수를 역산하는 방법(Fan and Nguyen 각주 24)이 유일한 경로다. V6 API에서는 Claude Code가 빠지고 AEI API는 직접 호출(1P)만 포함하므로, 매출에서 Claude Code와 클라우드 파트너 경유분을 빼야 한다. 이 두 몫은 공개된 값이 없어 가정으로 둘 수밖에 없다. (3) 따라서 LCE는 \textbf{단위당 가치}(데이터로 확정)와 \textbf{규모}(가정에 의존)를 나눠 보고하고, 규모는 격자형 민감도 표로 제시할 것을 권고한다. 외부 파라미터 원자료는 \texttt{data/raw/macro/scaling/}에 출처$\cdot$스냅샷과 함께 보관했다.
\end{abstract}

\tableofcontents

% ============================================================
\section{문제: 무엇이 빠져 있는가}
% ============================================================

\textcite{fan-2026-aggregate-gains-from-ai}의 LCE 식(식 1)은 다음과 같다.
\begin{equation*}
\text{LCE} = \sum_c \sum_s \sum_{t\in s} \frac{n_{c,t}}{N^*}\times C_{\text{total}} \times h_t \times \frac{W_{c,s}\times 12}{2{,}080}
\end{equation*}

report 10(6절)과 앞선 검토에서 확인했듯이, V6 원자료로 채울 수 있는 요소와 채울 수 없는 요소는 표~\ref{tab:avail}와 같다.

\begin{longtable}{L{3.4cm}L{5.6cm}L{5.6cm}}
\caption{LCE 요소별 V6 원자료 가용성}\label{tab:avail}\\
\toprule
요소 & Claude.ai (\texttt{aei\_claude\_ai\_2026-06-26.csv}) & 1P API (\texttt{aei\_1p\_api\_2026-06-26.csv}) \\
\midrule
\endhead
과업$\cdot$직업 구성 ($n_t/N$) & ○ \texttt{pct} (GLOBAL$\cdot$국가) & ○ \texttt{pct} (GLOBAL만) \\
절감시간 $h_t$ & ○ \texttt{human\_only\_time\_mean}(시간) $-$ \texttt{human\_with\_ai\_time\_mean}(분) & ○ 같은 변수 \\
임금 $W$ & 외부(ILO) & 외부. 국가 구분이 없어 단일 임금 벡터를 가정해야 함 \\
\textbf{규모 $C$} & × (\texttt{usage\_pct}=100뿐) & × \\
\bottomrule
\end{longtable}

V6 GLOBAL \texttt{overall} 행에서 단위당 절감시간은 이미 계산할 수 있다(2026년 5월 기준). Claude.ai는 대화당 약 4.06시간(4.73시간 $\to$ 40.1분), 1P API는 호출당 약 1.05시간(1.13시간 $\to$ 4.53분)이다. 따라서 이 노트의 과제는 \textbf{$C$를 어떻게 정할 것인가} 하나로 좁혀진다.

% ============================================================
\section{출발점: Fan and Nguyen(2026)의 규모 가정}
% ============================================================

원문 부록 B와 각주 24를 다시 정리하면 다음과 같다(자세한 인용은 \texttt{scaling\_parameters\_public.csv}의 \texttt{source\_tier=paper} 행).

\paragraph{Claude.ai (본문 LCE)}
\begin{enumerate}[label=(\arabic*)]
  \item 웹 MAU 1,890만 명(2026년 초, DemandSage가 Semrush$\cdot$SimilarWeb 추정치를 옮긴 수치)
  \item $\times$ DAU/MAU 50\%(업계 벤치마크 30--50\%의 상한, 벤치마크 출처는 없음)
  \item $\times$ 1인당 하루 3건(근거 없음) $\times$ 7일 $\to$ \textbf{주 2억 건}(범위 1.5억--2.5억)
  \item $\div$ 전문가용 AI 대화 중 Claude 점유율 25\%(범위 15--35\%) $\to$ \textbf{전체 AI 플랫폼 주 8억 건}($C_{\text{total}}$)
  \item 다른 웨이브는 SimilarWeb claude.ai 월 방문 수 비율로 늘리거나 줄임(2025.8 1.49억, 2025.11 1.76억, 2026.2 2.88억 회)
\end{enumerate}

\paragraph{1P API (각주 24, 공식 추정치에서는 제외)}
\begin{itemize}
  \item 연 호출 수 = run-rate 매출 140억 달러 $\times$ API 비중 75\% $\div$ 호출당 0.34달러 $\approx$ \textbf{연 310억 건}(주 약 6억 건)
  \item 물량의 절반을 미국 임금으로 환산 $\to$ 연 약 1.14조 달러(Anthropic API만, 점유율로 확장하지 않음)
\end{itemize}

두 가정 모두 LCE 수준에 정비례하며, 분포 지표(ACI)에는 영향을 주지 않는다.

% ============================================================
\section{V6 기간에 쓸 수 있는 외부 자료}
% ============================================================

\begin{longtable}{L{3.0cm}L{2.9cm}L{1.9cm}L{4.4cm}L{2.4cm}}
\caption{규모 추정용 외부 자료(2026-09-27 조사)}\label{tab:sources}\\
\toprule
변수 & 값 & 시점 & 출처 & 신뢰도 \\
\midrule
\endhead
claude.ai 월 방문 수 & 2.88억 & 2026.2 & Fan and Nguyen(2026) 부록 B (SimilarWeb) & 논문 \\
 & 6.14억 & 2026.3 & Similarweb 보고서, PPC Land 인용 & 2차 \\
 & \textbf{8.24억} & \textbf{2026.4} & 같음 & 2차 \\
 & \textbf{9.47억--9.69억 범위} & \textbf{2026.5--8} & 같음 & 2차 \\
 & 9.467억 (전월 대비 $-$0.61\%) & 2026.6 & DemandSage (Similarweb) & 2차 \\
Claude 앱 MAU & 5,510만 $\to$ 1.02억 $\to$ 1.20억 & 2026.3 $\to$ 5 $\to$ 6 & Similarweb App Intelligence, PPC Land 인용 & 2차 \\
Claude 웹+앱 MAU & 2.45억 & 2026년 중반 & DemandSage (Sensor Tower) & 2차 \\
생성형 AI 웹 방문 합계 & 월평균 95억 & 2025.6--2026.5 & Similarweb 블로그 & Similarweb 직접 \\
ChatGPT 웹 트래픽 점유율 & 약 76\% $\to$ 약 53\% & 2025.6 $\to$ 2026.5 & 같음 & Similarweb 직접 \\
Claude 웹 트래픽 점유율 & ``10\%에 근접'' & 2026년 하반기 & 같음 & Similarweb 직접(정성적) \\
Anthropic run-rate 매출 & 140억 달러 & 2026.2.12 & Anthropic 공식(Series G) & 공식 \\
 & 300억 달러 이상 & 2026.4.6 & Anthropic 공식(Google$\cdot$Broadcom 발표) & 공식 \\
 & 470억 달러 이상 & 2026.5 & Anthropic 공식(Series H, 5.28) & 공식 \\
Claude Code run-rate 매출 & 25억 달러 이상 & 2026.2.12 & Anthropic 공식(Series G) & 공식 \\
API 매출 비중 & 75--85\% & 2026 & Panto AI 등 제3자 & 미검증 \\
\bottomrule
\end{longtable}

자료의 성격에 대해 세 가지를 짚어 둔다.
\begin{itemize}
  \item \textbf{공식 수치는 매출뿐이다.} Anthropic은 이용자 수$\cdot$대화 수$\cdot$토큰 양을 공개하지 않는다. Claude Code 매출도 2월 이후 수치가 없다.
  \item \textbf{Similarweb 시계열 원자료는 유료다.} 위 방문 수는 대부분 2차 기사를 거친 값이고, 서로 다른 조회 시점(vintage)이 섞여 있다. 2026년 5월 방문 수는 점값으로 보고된 적이 없어, 6월 값과 전월 대비 변화율로 역산해야 한다($9.467\text{억}\div(1-0.0061)\approx 9.525\text{억}$).
  \item \textbf{앱 MAU는 출처별로 크게 다르다.} Fan and Nguyen이 인용한 2026년 2월 앱 MAU 1,250만 명(DemandSage)과 Similarweb의 3월 5,510만 명은 한 달 사이 4.4배 차이다. 두 값은 서로 다른 계열이므로 비율을 내면 안 된다.
\end{itemize}

% ============================================================
\section{Claude.ai 규모 추정 방법}
% ============================================================

\subsection{방법 A(기준): 웹 방문 수 비율로 늘리기 --- Fan and Nguyen 방식의 연장}

\begin{equation*}
C^{\text{Claude}}_m = 2\text{억} \times \frac{V_m}{V_{2026.2}}
\end{equation*}

\begin{longtable}{L{2.8cm}L{3.4cm}L{2.6cm}L{2.8cm}L{3.0cm}}
\caption{방법 A에 따른 V6 기간 추정치}\label{tab:methodA}\\
\toprule
월 & 방문 수 $V_m$ & $V_m/V_{2026.2}$ & Claude.ai 대화/주 & $C_{\text{total}}$/주 (점유율 25\%) \\
\midrule
\endhead
2026.2 (R5, 기준) & 2.88억 & 1.00 & 2.0억 & 8.0억 \\
2026.4 & 8.24억 & 2.86 & \textbf{5.7억} & 22.9억 \\
2026.5 & 약 9.53억 (범위 9.47억--9.69억) & 3.31 (3.29--3.36) & \textbf{6.6억} (6.6억--6.7억) & 26.5억 \\
\bottomrule
\end{longtable}

이 방법의 장점은 Fan and Nguyen의 R3--R5 추정치와 같은 방식이라 시계열로 이어진다는 점이다. 약점은 다음 네 가지다.
\begin{enumerate}[label=(\arabic*)]
  \item \textbf{대화/방문 비율이 고정되어 있다는 가정.} 2월에서 5월 사이에 앱 MAU가 급증했고(3월 5,510만 $\to$ 5월 1.02억), V6부터는 Cowork(데스크톱 앱)도 Claude.ai에 포함된다. 웹 방문만으로 늘리면 앱$\cdot$데스크톱 사용분이 빠져 \textbf{과소추정}할 가능성이 있다.
  \item \textbf{신규 이용자 효과.} 3--4월 유입자는 방문당 대화 수가 기존 이용자와 다를 수 있다.
  \item \textbf{vintage 불일치.} 2월 값(논문)과 4--5월 값(2차 기사)이 서로 다른 시점에 조회되었다.
  \item \textbf{2월 앵커 자체의 불확실성.} 주 2억 건은 근거가 약한 가정 두 개(DAU/MAU 50\%, 하루 3건)에 의존한다. 1.5억--2.5억 범위를 함께 들고 가야 한다.
\end{enumerate}

\subsection{방법 B(보완): 웹$\cdot$앱 결합 지수}

\begin{equation*}
C^{\text{Claude}}_m = 2\text{억} \times \left[\omega \frac{V_m}{V_{2026.2}} + (1-\omega)\frac{A_m}{A_{2026.2}}\right]
\end{equation*}

$A_m$은 앱 MAU, $\omega$는 2월 기준 웹 사용 비중이다. 앱 MAU는 \textbf{같은 계열(Similarweb App Intelligence)의 2월 값}이 있어야 계산할 수 있으므로, Similarweb 원자료를 받은 뒤에 적용한다. $\omega$는 관측되지 않으므로 0.5--0.9 범위로 민감도를 본다. 앱 성장률이 웹보다 높았기 때문에, 이 방법은 방법 A보다 큰 값을 준다.

\subsection{전체 AI 플랫폼으로 확장할 때($C_{\text{total}}$): 점유율 25\%를 고정하지 않는다}

Fan and Nguyen의 25\%는 ``전문가용 AI 대화 중 Claude 점유율''이다. 그런데 2월에서 5월 사이 Claude 트래픽이 3배 넘게 늘어난 반면, 생성형 AI 전체 웹 트래픽은 그만큼 늘지 않았다(ChatGPT 점유율 76\% $\to$ 53\%, Claude는 10\%에 근접). 따라서 점유율을 25\%로 고정하면 \textbf{$C_{\text{total}}$이 과대 추정}된다. 대안은 두 가지다.
\begin{itemize}
  \item \textbf{방법 C1(권장): 카테고리 트래픽으로 직접 늘리기.} $C_{\text{total},m} = 8\text{억}\times G_m/G_{2026.2}$. 여기서 $G$는 생성형 AI 웹 방문 합계(claude.ai, chatgpt.com, gemini.google.com, copilot, perplexity 등)다. Claude의 점유율 가정이 필요 없어진다. Similarweb 원자료가 필요하다.
  \item \textbf{방법 C2: 점유율을 웹 점유율 변화에 맞춰 조정하기.} $\text{share}_m = 25\%\times(\text{Claude 웹 점유율}_m/\text{Claude 웹 점유율}_{2026.2})$. 전문가용 점유율과 웹 점유율이 같은 비율로 움직였다는 가정이 들어간다.
\end{itemize}
어느 쪽이든 \textbf{Claude.ai만의 LCE}(점유율 가정 없음)와 \textbf{전체 플랫폼 LCE}(점유율 또는 카테고리 가정)를 따로 보고하는 편이 가정의 층위를 드러내는 데 유리하다.

% ============================================================
\section{1P API 규모 추정 방법}
% ============================================================

\subsection{추정식}

Fan and Nguyen 각주 24의 식을 V6의 데이터 범위에 맞게 두 항을 더해 확장한다.
\begin{equation*}
Q^{\text{API}}_m\ (\text{연간 호출 수}) = \frac{R_m \times s_{\text{API}} \times (1-\kappa_{\text{CC}}-\kappa_{\text{cloud}})}{p}
\end{equation*}

\begin{longtable}{L{1.4cm}L{3.8cm}L{3.2cm}L{6.2cm}}
\caption{API 추정식의 파라미터}\label{tab:apiparam}\\
\toprule
기호 & 의미 & 값/범위 & 근거 \\
\midrule
\endhead
$R_m$ & 전사 run-rate 매출 & 4월 $\geq$300억, 5월 $\geq$470억 달러 & Anthropic 공식 \\
$s_{\text{API}}$ & 매출 중 API(사용량 과금) 비중 & 0.75 (0.70--0.85) & Fan and Nguyen 0.75, 제3자 75--85\% \\
$\kappa_{\text{CC}}$ & API 매출 중 Claude Code 몫 & 0.15--0.30 (가정) & V6 API는 Claude Code 제외. 2월 Claude Code는 전사 매출의 18\% 이상(25억/140억)이었지만 구독 과금분이 섞여 있어 API 내 몫은 불명 \\
$\kappa_{\text{cloud}}$ & API 매출 중 Bedrock$\cdot$Vertex 등 클라우드 경유 몫 & 0.2--0.4 (가정) & AEI API는 1P만 포함. 공개 수치 없음 \\
$p$ & 호출당 매출 & 0.34달러 (0.2--0.5) & Fan and Nguyen. 도출 근거 없음. 모델 구성$\cdot$가격 변화 반영 필요 \\
\bottomrule
\end{longtable}

주간 값은 $Q/52$로 구한다.

\subsection{예시 계산}

\begin{longtable}{L{6.4cm}L{2.6cm}L{2.6cm}L{2.6cm}}
\caption{1P API 주간 호출 수 예시}\label{tab:apiex}\\
\toprule
사양 & 2026.2 (참고) & 2026.4 & 2026.5 \\
\midrule
\endhead
Fan 방식 그대로($\kappa=0$, $p=0.34$, $s=0.75$) & 5.9억 & 12.7억 & 19.9억 \\
제외분 반영($\kappa_{\text{CC}}=0.2$, $\kappa_{\text{cloud}}=0.3$) & 3.0억 & 6.4억 & 10.0억 \\
\bottomrule
\end{longtable}

첫 줄은 Claude Code와 클라우드 경유분이 섞인 값이라 V6 1P API의 \textbf{상한}으로 읽어야 한다. 두 줄의 차이만 봐도 알 수 있듯이, 이 경로에서는 공개되지 않은 몫($\kappa$)이 결과를 두 배 가까이 바꾼다.

\subsection{주의점}

\begin{enumerate}[label=(\arabic*)]
  \item \textbf{run-rate는 시점 값이다.} 300억 달러는 4월 6일, 470억 달러는 ``5월 초''의 값이며 모두 하한(``surpassed'', ``crossed'')이다. 월평균으로 쓰려면 두 점 사이를 로그 선형으로 보간해 월 중순 값을 만드는 것을 대안 사양으로 둔다.
  \item \textbf{매출과 호출 수의 관계가 불안정하다.} 모델 구성(Opus 비중)과 요금, 캐싱 할인에 따라 호출당 매출 $p$가 달라진다. V6 원자료에는 토큰$\cdot$비용 정보가 없어 $p$를 데이터로 보정할 수 없다.
  \item \textbf{호출은 대화가 아니다.} V6 API 레코드는 입력--출력 한 쌍이다. 파이프라인에서 여러 번 호출해 과업 하나를 끝내면, 호출마다 절감시간($h_t$)을 더할 경우 \textbf{이중으로 세게 된다}. 호출당 절감시간(약 1.05시간)이 대화당(약 4.06시간)보다 훨씬 작은 것은 이 점을 일부 반영하지만, 과대 추정 위험은 여전히 남는다.
  \item \textbf{점유율로 확장하지 않는다.} API LCE는 Anthropic 1P API만의 값이다. Claude.ai의 전체 플랫폼 LCE와 더하거나 수준을 비교하면 안 된다.
\end{enumerate}

% ============================================================
\section{두 플랫폼의 상대 규모: report 10과의 연결}
% ============================================================

report 10(6.2절)은 두 플랫폼 사용량을 합산하려면 가중치 $w$(Claude.ai 몫)를 외생적으로 정해야 한다고 결론 내렸다. 위 추정은 이 $w$에 대한 외부 근거를 준다. 2026년 5월 기준으로 방법 A의 Claude.ai 주 6.6억 건과 표~\ref{tab:apiex}의 API 주 10억--20억 건을 쓰면, $\boldsymbol{w\approx 0.25\text{--}0.40}$이다. 표본 동일 가중($w=0.5$, \textcite{massenkoff-2026-labor-market-impacts-of-ai}의 암묵적 가정)보다 API 쪽 비중이 크다. 다만 대화 1건과 호출 1건이 같은 단위가 아니라는 점(5.3절 3항)을 함께 적어야 한다.

% ============================================================
\section{권고: 보고 방식}
% ============================================================

\begin{enumerate}[label=(\arabic*)]
  \item \textbf{단위당 가치와 규모를 분리한다.} 대화(호출)당 LCE, 즉 $\sum_s \text{pct}_s \times h_s \times w_s$는 데이터로 확정되는 값이다. 여기에 규모를 곱한 총량은 가정에 의존하는 값으로 구분해 보고한다.
  \item \textbf{민감도 격자를 제시한다.}
  \begin{itemize}
    \item Claude.ai: 2월 앵커(1.5억/2.0억/2.5억) $\times$ 확장 방법(A/B) $\times$ (전체 플랫폼일 때) 점유율 또는 카테고리 확장(C1/C2)
    \item API: $R$(공식 하한/보간) $\times$ $s_{\text{API}}$(0.70/0.75/0.85) $\times$ $\kappa_{\text{CC}}$(0.15/0.30) $\times$ $\kappa_{\text{cloud}}$(0.2/0.4) $\times$ $p$(0.2/0.34/0.5)
  \end{itemize}
  \item \textbf{월별로 따로 계산한다.} 4월과 5월 사이에 매출이 1.5배로 늘었으므로 두 달을 평균하지 않는다.
  \item \textbf{파라미터는 raw 파일에서 읽는다.} \texttt{data/raw/macro/scaling/scaling\_parameters\_public.csv}를 do 파일에서 불러오고, 5월 방문 수 역산 같은 파생값도 do 파일에서 계산한다. 수치를 본문에 직접 쓰지 않는다.
\end{enumerate}

% ============================================================
\section{후속 작업}
% ============================================================

\begin{enumerate}[label=(\arabic*)]
  \item \textbf{Similarweb 원자료 확보}: 월 방문 수(claude.ai와 주요 경쟁 서비스, 2025.1--최신)와 앱 MAU를 같은 날 한 번에 내려받아 \texttt{data/raw/macro/scaling/}에 둔다. 목록은 해당 폴더 \texttt{README.md} 참고. 이것이 확보되어야 방법 B$\cdot$C1을 계산할 수 있고 vintage 문제도 해소된다.
  \item \textbf{do 파일}: \texttt{code/01\_import\_scaling.do}(파라미터 불러오기$\cdot$파생값 계산 $\to$ \texttt{data/proc/macro/scaling/}), \texttt{code/04\_analysis\_lce\_global.do}(V6 GLOBAL 단위당 가치 $\times$ 규모 격자 $\to$ \texttt{results/table/}).
  \item \textbf{API 임금 가정}: 국가 구분이 없는 API에 쓸 임금 벡터(미국 OES, 미국 50\% + Claude.ai 국가 구성 가중 ILO 50\%, Claude.ai GLOBAL 사용 가중 ILO)를 정한다.
\end{enumerate}

% ============================================================
\section*{웹 자료}
% ============================================================

접속일은 모두 2026-09-27이다. URL, 스냅샷 파일은 \texttt{data/raw/macro/scaling/scaling\_parameters\_public.csv}에 행 단위로 기록했다.
\begin{itemize}
  \item Anthropic (2026.2.12). Series G 발표. \url{https://www.anthropic.com/news/anthropic-raises-30-billion-series-g-funding-380-billion-post-money-valuation}
  \item Anthropic (2026.4.6). Google$\cdot$Broadcom 파트너십 확대 발표. \url{https://www.anthropic.com/news/google-broadcom-partnership-compute}
  \item Anthropic (2026.5.28). Series H 발표. \url{https://www.anthropic.com/news/series-h}
  \item Similarweb (2026). ``AI Search Stats 2026''(2026.7.29 게시, 9.17 수정). \url{https://aisearch.similarweb.com/blog/gen-ai-stats/}
  \item PPC Land (2026.9). ``Claude gains 540\% in visits with 0.46\% from paid traffic, Similarweb finds.'' \url{https://ppc.land/claude-gains-540-in-visits-with-0-46-from-paid-traffic-similarweb-finds/}
  \item DemandSage. ``Claude AI Statistics (2026).'' \url{https://www.demandsage.com/claude-ai-statistics/}
\end{itemize}

\printbibliography[title={참고문헌}]

\end{document}
